\section{Discussion and conclusion}
\label{sec:discussion}

The central technical result of this work is that two etch-rate
anchors at a single $(E/N, x_{\mathrm{Ar}})$ operating condition
constrain exactly one scalar combination of the depletion and
rate-coefficient parameters of a reduced \sfsix{} ICP model. The
identifiable invariant is the product
$\beta\,\kdiss\,\nel(P_{\mathrm{high}})$, with the electron density
computed self-consistently from the power balance and the overall
magnitude absorbed by the etch yield $\varepsilon$. This
identifiability structure, derived analytically in the isothermal
limit and verified numerically in the heated model
(Section~\ref{sec:identifiability}), explains why a
Maxwellian-calibrated reduced model can simultaneously reproduce
the calibration anchors to machine precision and produce a fitted
depletion timescale that cannot be interpreted as a physical
residence time: the Maxwellian rate-coefficient convolution
overestimates $\kdiss$ in attachment-dominated \sfsix{} at the
operating electron temperature, and the error is absorbed into
the single fitted parameter that participates in the identifiable
product alongside $\kdiss$. Replacing the Maxwellian rates with
two-term Boltzmann rates from BOLSIG$^+$ shifts the fitted
depletion parameter by the factor predicted analytically from the
identifiability invariant (Section~\ref{sec:boltzmann}), into
order-of-magnitude agreement with a nominal geometric
residence-time range. At the calibration operating condition, the
Maxwellian- and Boltzmann-calibrated models are observationally
indistinguishable for any set of absorbed-power measurements; off
the calibration axis, their predictions diverge by up to $85\,\%$
along the reduced-field axis and $47\,\%$ along the Ar-dilution
axis, with a well-defined observability maximum at intermediate Ar
fraction $\approx\!30\,\%$ (Section~\ref{sec:observability}).

The methodological contribution is a three-stage workflow that we
propose as a minimum standard for assigning physical
interpretation to fitted parameters in reduced
electronegative-plasma models. The first stage is an identifiability
analysis of the calibration problem, carried out analytically where
possible and numerically where the model structure resists
closed-form manipulation; its goal is to determine which scalar
combinations of the model parameters are actually constrained by
the data. The second stage is a rate-source validation, in which
the sensitivity of the fitted parameters to the rate-coefficient
computation method is quantified by recalibrating against the same
data with an alternative rate source; a parameter whose fitted
value shifts between rate sources by an amount predicted by the
identifiability invariant is not a physical quantity but a proxy
for the rate-source assumption. The third stage is an off-axis
observability characterization that maps operating-condition
configurations at which competing rate sources predict measurably
different outcomes, so that experimental campaigns can be targeted
at observationally useful perturbations. None of the three stages
is novel in isolation; the contribution is the claim that all
three are necessary and that omitting any one produces fitted
parameters whose physical interpretation is at best uncontrolled.

Several limitations of the present study bound the generality of
the conclusions, and we state them explicitly. First, the work
contains no independent experimental validation. The
\citet{mansano1999} and \citet{panduranga2019} etch-rate anchors
are drawn from different reactors, and the agreement between the
Boltzmann-calibrated depletion timescale and a nominal geometric
residence-time range is a self-consistency check against a
textbook mass-balance estimate, not a direct measurement against
an observed residence time in the calibrated reactor. The paper
should therefore be read as a methodology paper with a worked
example: the identifiability theorem of
Section~\ref{sec:identifiability} and the rate-source-swap test of
Section~\ref{sec:boltzmann} are methodological results whose
validity does not depend on the specific numerical value of the
fitted $\beta$, while the physical interpretability claim is a
consistency check that the experimental follow-up recommended in
Section~\ref{sec:observability} remains necessary to close.

Second, we treat the two-term Boltzmann solution from BOLSIG$^+$ as
the rate-source comparison target, but this is itself an
approximation of the electron kinetic equation, valid to lowest
order in the angular expansion and specific to the chosen
cross-section library \citep{biagi}. The claim of this work is
specifically that the Maxwellian-to-Boltzmann shift captures the
dominant EEDF bias in the attachment-dominated regime; residual
biases from higher-order EEDF effects (multi-term Boltzmann,
anisotropy, particle-in-cell corrections) would propagate through
the same identifiability invariant \eqref{eq:invariant-full} and
shift the fitted $\beta$ by an additional factor predictable from
equation~\eqref{eq:beta-shift}. The identifiability analysis is
therefore agnostic to the specific rate source used beyond the
Maxwellian; what it requires is a rate source that is closer to
the true EEDF than a Maxwellian at the same mean energy.

Third, the reduced model itself makes structural assumptions that
our methodology does not probe: one-dimensional radial diffusion
with Robin wall boundary conditions, a single-channel power
balance retaining only ionization and attachment losses, a
one-parameter feedstock-depletion correction, and a linear
gas-heating coefficient fitted to the calibration data. Errors in
any of these assumptions are orthogonal to the EEDF error
identified here and would require a separate analysis.

Fourth, the two literature anchors used to calibrate the model are
from chambers operated at different pressures
(\citet{mansano1999} at $10\,$mTorr with a $150\,$W cathode bias,
\citet{panduranga2019} at $30\,$mTorr with zero bias). The
identifiability theorem of
Section~\ref{sec:identifiability} is pressure-independent and the
rate-source-swap prediction test is unaffected, but the absolute
numerical value of $\beta$ (and hence its order-of-magnitude
comparison to geometric residence time in
Section~\ref{sec:boltzmann}) is specific to the $10\,$mTorr
nominal operating condition into which we absorbed the systematic.
A future validation experiment in a single chamber at a single
pressure, with at least three power levels and optional Ar
dilution, would remove this systematic and is the natural next
step. The observability analysis of
Section~\ref{sec:observability} identifies Ar dilution at
$x_{\mathrm{Ar}} \approx 0.3$ and $200\,$W as the operating point
with highest Maxwell--Boltzmann discrimination and therefore the
most informative single experiment.

The methodology generalizes to other electronegative-plasma reduced
models in which the source term for a reactive neutral is governed
by electron-impact dissociation of a strongly attaching feedstock.
The attachment-induced depletion of the EEDF tail responsible for
the large Maxwellian bias in \sfsix{} rate coefficients appears
with different magnitudes but identical structure in Cl$_2$,
BCl$_3$, CF$_4$, and NF$_3$ plasmas, and the identifiability
analysis depends only on the functional form of the depletion
correction and not on the specific gas. A reduced model calibrated
against etch-rate data for any of these gases with Maxwellian rate
coefficients should therefore be expected to produce a fitted
depletion timescale whose literal value inherits a multiplicative
bias of the same form as the \sfsix{} case. The broader message
for reduced-order plasma modeling is that parameter identifiability
is a prerequisite for physical interpretation, not a diagnostic to
run afterwards.

This paper establishes three results. The two-anchor depletion
calibration of a reduced \sfsix{} ICP model is structurally
degenerate in the rate-coefficient source, and the fitted
depletion parameter acquires a multiplicative EEDF-dependent bias
that cannot be diagnosed from the calibration data alone. The
Maxwellian-EEDF assumption is responsible for the discrepancy
between the fitted depletion timescale and a plausible residence
time for the target reactor class, as demonstrated by a controlled
rate-source swap whose outcome is quantitatively predicted by the
identifiability invariant to better than $0.5\,\%$. The calibration
degeneracy is structural at the anchor operating condition but
opens up operationally at off-anchor configurations, with a
practical experimental handle (Ar dilution at $\sim\!30\,\%$) that
can distinguish the two competing EEDF hypotheses in a single
additional measurement. The significance of the work is not in the
updated numerical value of the fitted parameter, but in the
demonstration that reduced-order plasma models require
identifiability analysis and rate-source validation before their
fitted parameters can carry physical interpretation, and that both
steps are tractable with the same computational apparatus used for
the calibration itself.
